
% Integration-ready chapter fragment.
\chapter{Shifted RTT duality, orthogonal coideal descent, and the rank-one Kleinian test case}
\label{chap:shifted-rtt-orthogonal-coideals}

\section{Four problems that completion conflates}
\label{sec:purpose-status-discipline}

The Yangian/shifted-Yangian story conflates four issues that must be separated
before any theorem can be stated cleanly.

\begin{enumerate}[label=\textup{(\roman*)}]
\item the \emph{ambient category} in which bar/cobar inversion holds;
\item the \emph{local RTT kernel}, which is finite-stage and quadratic;
\item the \emph{boundary quotient}, where genuine shifted/truncated data enter;
\item the \emph{rank-one test case}, where the orthogonal-coideal mechanism can
 be computed explicitly and matched to a Kleinian-type algebra.
\end{enumerate}

The first point is governed by completed/pro-nilpotent chiral Koszul duality and
factorization homology \cite{FG12,AF15,AF20}. The shifted-Yangian side is
motivated by the classical shifted Yangians of Brundan--Kleshchev and by the
Coulomb-branch picture of Braverman--Finkelberg--Nakajima
\cite{BrunKlesh06,BFN18}; the dg-shifted line-operator target is the recent
construction of Dimofte--Niu--Py \cite{DNP25}. Within this chapter, every
statement is intended to be read literally at the level at which it is stated:
construction versus inversion, local kernel versus boundary quotient, finite
stage versus completed limit, and theorem versus conjecture are kept sharply
separate.

\begin{definition}[Three layers of the ordered theory]
\label{def:three-layers-ordered-theory-appendix}
The ordered/factorization Yangian story decomposes into three layers.
\begin{enumerate}[label=\textup{(\alph*)}]
\item The \emph{local RTT kernel}: a finite-stage quadratic $E_1$-presentation
 controlled by the rational $R$-matrix.
\item The \emph{completed envelope}: the weightwise stabilized bar/cobar or
 Maurer--Cartan object in the completed/pro-nilpotent setting.
\item The \emph{boundary quotient}: the extra datum cutting the unshifted local
 kernel down to a shifted/truncated or Coulomb-branch-type object.
\end{enumerate}
A foundational correction is that the shift belongs primarily to layer
\textup{(c)}, not to layer \textup{(a)} alone.
\end{definition}

\begin{remark}[Why the distinction matters]
\label{rem:why-the-distinction-matters}
Layer \textup{(a)} is where one can write explicit formulas and prove exact
adjointness theorems. Layer \textup{(b)} is where infinite-generator towers stop
being inverse-limit pathologies and become weightwise stabilization problems.
Layer \textup{(c)} is where genuine shifted/truncated geometry enters. If one
tries to force the boundary data into the local quadratic kernel alone, one
usually either gets a tautological change of basis or a quotient that kills the
root sector. The rank-one computation below shows this explicitly.
\end{remark}

\section{Weightwise stabilized completed bar/cobar envelopes}
\label{sec:weightwise-stabilized}

\begin{definition}[Positive weight-graded complete $A_\infty$-algebra]
\label{def:positive-weight-complete-ainfty}
A \emph{positive weight-graded complete $A_\infty$-algebra} is a complete
augmented $A_\infty$-algebra
\[
A=\prod_{w\ge 0}A_w
\]
with structure maps
\[
m_n\colon A_{w_1}\otimes\cdots\otimes A_{w_n}\longrightarrow A_{w_1+\cdots+w_n}
\qquad (n\ge 1)
\]
preserving total weight, and, in the curved case, curvature
\[
m_0\in \prod_{w\ge 1} A_w.
\]
Write
\[
A_{\le N}:=\bigoplus_{0\le w\le N}A_w.
\]
\end{definition}

\begin{lemma}[Weightwise stabilization of the bar complex \ClaimStatusProvedHere]
\label{lem:weightwise-stabilization-bar-complex}
For every bar weight $w$, the weight-$w$ summand of the reduced bar complex
depends only on the truncation $A_{\le w}$:
\[
\overline{B}(A)_w\cong \overline{B}(A_{\le N})_w\cong \overline{B}(A_{\le w})_w
\qquad (N\ge w).
\]
In particular, every chain group and differential on bar weight $w$ stabilizes
once the truncation level reaches $w$.
\end{lemma}

\begin{proof}
A bar monomial of total weight $w$ has the form
\[
sa_1|\cdots|sa_r,
\qquad \sum_j \mathrm{wt}(a_j)=w.
\]
Hence every $a_j$ has weight at most $w$. The bar differential is built from
the $m_n$, and each $m_n$ preserves total weight. Therefore the weight-$w$
piece of the differential only uses the structure maps on weights $\le w$.
Truncating above weight $w$ changes nothing.
\end{proof}

\begin{definition}[Stabilized completed bar coalgebra]
\label{def:stabilized-completed-bar-coalgebra}
The \emph{stabilized completed bar coalgebra} of a positive weight-graded
complete $A_\infty$-algebra $A$ is
\[
\mathbb{B}_{\mathrm{st}}(A)_w:=\overline{B}(A_{\le w})_w,
\qquad
\widehat{\mathbb{B}}_{\mathrm{st}}(A):=\prod_{w\ge 0}\mathbb{B}_{\mathrm{st}}(A)_w.
\]
By Lemma~\ref{lem:weightwise-stabilization-bar-complex}, this is independent of
any auxiliary truncation level once the weight is fixed.
\end{definition}

\begin{theorem}[Stabilized completed bar/cobar recovery \ClaimStatusProvedHere]
\label{thm:stabilized-completed-recovery}
Let $A$ be a positive weight-graded complete $A_\infty$-algebra such that every
finite truncation $A_{\le N}$ lies in the finite-type bar/cobar regime and the
finite-stage counits
\[
\epsilon_N\colon \Omega\overline{B}(A_{\le N})\xrightarrow{\sim} A_{\le N}
\]
are quasi-isomorphisms. Then the completed counit
\[
\widehat\epsilon\colon \Omega\widehat{\mathbb{B}}_{\mathrm{st}}(A)
\longrightarrow A
\]
is a quasi-isomorphism in the completed graded category.
\end{theorem}

\begin{proof}
Fix a weight $w$. By construction,
\[
\Omega\widehat{\mathbb{B}}_{\mathrm{st}}(A)_w
\cong
\Omega\overline{B}(A_{\le w})_w,
\]
because every tensor word of total weight $w$ only uses coalgebra summands of
weight at most $w$. Under this identification, the weight-$w$ piece of the
completed counit is exactly the weight-$w$ piece of
\[
\epsilon_w\colon \Omega\overline{B}(A_{\le w})\xrightarrow{\sim} A_{\le w}.
\]
Hence it is a quasi-isomorphism. Cohomology in the completed graded category is
computed weightwise, so the product over all weights is again a
quasi-isomorphism.
\end{proof}

\begin{corollary}[Automatic tower convergence in positive degree \ClaimStatusProvedHere]
\label{cor:automatic-tower-convergence-positive-degree}
For an honest positive conformal/RTT grading, the infinite-generator completion
problem is not an inverse-limit problem. It is a finite-stage problem followed
by weightwise stabilization.
\end{corollary}

\begin{theorem}[Weightwise Maurer--Cartan descent \ClaimStatusProvedHere]
\label{thm:weightwise-maurer-cartan-descent}
Let
\[
L=\prod_{w\ge 1}L_w
\]
be a complete positively weight-graded $L_\infty$-algebra whose structure maps
preserve total weight. Then the natural map
\[
\operatorname{MC}(L)\longrightarrow \varprojlim_N \operatorname{MC}(L_{\le N})
\]
is a bijection.
\end{theorem}

\begin{proof}
Write $\alpha=\sum_{w\ge 1}\alpha_w\in L^1$. The weight-$m$ part of the
Maurer--Cartan equation
\[
\sum_{n\ge 1}\frac{1}{n!}\ell_n(\alpha,\dots,\alpha)=0
\]
is
\[
\sum_{n=1}^{m}\frac{1}{n!}
\sum_{w_1+\cdots+w_n=m}
\ell_n(\alpha_{w_1},\dots,\alpha_{w_n})=0.
\]
This depends only on $\alpha_1,\dots,\alpha_m$, hence is exactly the
Maurer--Cartan equation in the truncation $L_{\le m}$. Therefore $\alpha$ is a
Maurer--Cartan element of $L$ if and only if each truncated piece
$\alpha_{\le N}$ is Maurer--Cartan in $L_{\le N}$.
\end{proof}

\begin{proposition}[Module/comodule passage for stabilized towers \ClaimStatusProvedHere]
\label{prop:module-comodule-stabilized}
Assume, in addition to the hypotheses of
Theorem~\ref{thm:stabilized-completed-recovery}, that each finite truncation
admits a bar-equivalence of module categories compatible with truncation:
\[
\overline{B}_{A_{\le N}}\colon \mathrm{Mod}^{\mathrm{pn}}(A_{\le N})
\xrightarrow{\sim}
\mathrm{Comod}^{\mathrm{conil}}(\overline{B}(A_{\le N})).
\]
Then passing weightwise to the stabilized completion produces an equivalence
\[
\mathrm{Mod}^{\mathrm{comp}}(A)
\xrightarrow{\sim}
\mathrm{Comod}^{\mathrm{cont}}(\widehat{\mathbb{B}}_{\mathrm{st}}(A)).
\]
\end{proposition}

\begin{proof}
The weight-$w$ sectors stabilize exactly as in
Lemma~\ref{lem:weightwise-stabilization-bar-complex}. Therefore the inverse
system of finite-stage module/comodule equivalences becomes a degreewise
constant system in every fixed weight. Passing to the product over weights
preserves the finite-stage equivalences.
\end{proof}

\section{An explicit bar-duality map and the Lagrangian envelope}
\label{sec:explicit-bar-duality-lagrangian-envelope}

\begin{definition}[Finite bar-duality datum]
\label{def:finite-bar-duality-datum}
A \emph{finite bar-duality datum} for a pair of augmented positive weight-graded
$A_\infty$-algebras $(A,A^!)$ consists of:
\begin{enumerate}[label=\textup{(\roman*)}]
\item finite-dimensional weight slices $A_w$ and $A^!_w$;
\item perfect pairings
 \[
 \langle -, -\rangle_w\colon s\overline{A}_w\otimes s\overline{A^!}_w\to \mathbf{k}[-1];
 \]
\item induced perfect pairings on reduced bar slices
 \[
 \langle -, -\rangle_{\overline{B},w}\colon \overline{B}(A)_w\otimes \overline{B}(A^!)_w\to \mathbf{k};
 \]
\item adjointness of the bar differentials:
 \[
 \langle D_Ax,y\rangle_{\overline{B},w}
 +(-1)^{|x|}\langle x,D_{A^!}y\rangle_{\overline{B},w}=0.
 \]
\end{enumerate}
Set
\[
\mathrm{Def}^{\mathrm{wt}}(A):=\mathrm{CoDer}^{\mathrm{wt}}(\overline{B}(A))[1],
\qquad
\mathrm{Def}^{\mathrm{wt}}(A^!):=\mathrm{CoDer}^{\mathrm{wt}}(\overline{B}(A^!))[1].
\]
\end{definition}

\begin{theorem}[Trace-theoretic bar duality \ClaimStatusProvedHere]
\label{thm:trace-theoretic-bar-duality}
From a finite bar-duality datum one obtains a canonical degree $-1$ chain
isomorphism
\[
\Theta_A\colon \mathrm{Def}^{\mathrm{wt}}(A^!)
\xrightarrow{\sim}
\mathrm{Def}^{\mathrm{wt}}(A)^\vee[-1].
\]
On weight $w$, if
\[
g_w\colon \overline{B}(A^!)_w\to s\overline{A^!}_w,
\qquad
f_w\colon \overline{B}(A)_w\to s\overline{A}_w,
\]
then one defines
\[
\Theta_{A,w}(g_w)(f_w):=\operatorname{Tr}(f_w\circ g_w^{\sharp}),
\]
where the bar pairing identifies
\[
g_w^{\sharp}\colon s\overline{A}_w\longrightarrow \overline{B}(A)_w[-1].
\]
\end{theorem}

\begin{proof}
Because the bar pairing is perfect on each weight slice,
\[
\overline{B}(A^!)_w\cong \overline{B}(A)_w^\vee,
\qquad
s\overline{A^!}_w\cong (s\overline{A}_w)^\vee[-1].
\]
Thus every $g_w$ dualizes to a map
\[
g_w^{\sharp}\colon s\overline{A}_w\to \overline{B}(A)_w[-1],
\]
and therefore $f_w\circ g_w^{\sharp}$ is an endomorphism of $s\overline{A}_w$ of
degree $-1$. The graded trace gives a perfect bilinear pairing between the two
Hom-spaces. Summing over all weights yields the desired map $\Theta_A$.

To check compatibility with the differentials, write the coderivation
differentials as graded commutators with the bar differentials:
\[
df=[D_A,f],
\qquad
dg=[D_{A^!},g].
\]
Adjointness of $D_A$ and $D_{A^!}$ implies
\[
(dg_w)^{\sharp}=-[D_A,g_w^{\sharp}].
\]
Hence
\[
\Theta_{A,w}(dg_w)(f_w)
+(-1)^{|g_w|+1}\Theta_{A,w}(g_w)(df_w)
=\operatorname{Tr}([D_A,f_wg_w^{\sharp}])=0,
\]
since the graded trace of a graded commutator vanishes. Therefore
$\Theta_A$ is a chain map.
\end{proof}

\begin{corollary}[The Lagrangian envelope \ClaimStatusProvedHere]
\label{cor:lagrangian-envelope}
Set
\[
\mathbb{T}_A:=\mathrm{Def}^{\mathrm{wt}}(A)\oplus \mathrm{Def}^{\mathrm{wt}}(A^!)[1].
\]
The bar-duality map of Theorem~\ref{thm:trace-theoretic-bar-duality} induces a
canonical $(-1)$-shifted symplectic pairing on $\mathbb{T}_A$, and the two
summands
\[
\mathrm{Def}^{\mathrm{wt}}(A)\subset \mathbb{T}_A,
\qquad
\mathrm{Def}^{\mathrm{wt}}(A^!)[1]\subset \mathbb{T}_A
\]
are complementary Lagrangians.
\end{corollary}

\begin{proof}
Define
\[
\omega\big((x,\xi),(y,\eta)\big)
:=\Theta_A(\eta)(x)-(-1)^{(|x|-1)(|y|-1)}\Theta_A(\xi)(y).
\]
This pairs only cross-terms, so each summand is isotropic. Perfectness of
$\Theta_A$ implies nondegeneracy. Because $\Theta_A$ is a chain map, $\omega$
is closed. Since $\mathbb{T}_A$ is the direct sum of the two isotropic
summands, each is maximal isotropic.
\end{proof}

\section{Meromorphic fiber functors and fusion reconstruction}
\label{sec:meromorphic-fiber-functors-fusion}

The monoidal/fusion problem can be reformulated entirely in terms of an exact
meromorphic fiber functor on a compact generating family.

\begin{lemma}[Generator extension of exact monoidal functors \ClaimStatusProvedHere]
\label{lem:generator-extension-exact-monoidal}
Let $F\colon \mathcal{C}\to \mathcal{D}$ be an exact colimit-preserving functor
between compactly generated stable monoidal categories. Let
$\mathcal{G}\subset \mathcal{C}^c$ be a set of compact generators. Assume:
\begin{enumerate}[label=\textup{(\roman*)}]
\item $F$ is fully faithful on $\mathcal{G}$;
\item the essential image $F(\mathcal{G})$ compactly generates $\mathcal{D}$;
\item for every $X,Y\in\mathcal{G}$, the comparison
 \[
 F(X)\otimes F(Y)\longrightarrow F(X\otimes Y)
 \]
 is an equivalence.
\end{enumerate}
Then $F$ is a monoidal equivalence.
\end{lemma}

\begin{proof}
Fix $Y\in\mathcal{C}$. The full subcategory of $X\in\mathcal{C}$ for which the
comparison $F(X)\otimes F(Y)\to F(X\otimes Y)$ is an equivalence is localizing
in $X$. By assumption it contains $\mathcal{G}$; since $\mathcal{G}$ generates,
it is all of $\mathcal{C}$. Repeating the argument in the $Y$-variable shows
that $F$ is monoidal on all objects.

Similarly, the full subcategory of $X$ for which
\[
\operatorname{Map}_{\mathcal{C}}(G,X)
\longrightarrow
\operatorname{Map}_{\mathcal{D}}(F(G),F(X))
\]
is an equivalence for every $G\in\mathcal{G}$ is localizing and contains the
generators, so $F$ is fully faithful. Its image is a localizing monoidal
subcategory containing the compact generators $F(\mathcal{G})$ of
$\mathcal{D}$, hence is all of $\mathcal{D}$.
\end{proof}

\begin{theorem}[Meromorphic Tannakian reconstruction \ClaimStatusProvedHere]
\label{thm:meromorphic-tannakian-reconstruction}
Let $\mathcal{L}$ be a stable $A_\infty$-category equipped with meromorphic
tensor products $\otimes_z$ for $z\in D^\times$. Let $B$ be a complete
$A_\infty$-algebra and suppose we are given exact functors
\[
F_z\colon \mathcal{L}\to \mathrm{Mod}^{\mathrm{comp}}(B)
\qquad (z\in D^\times\cup\{0\}).
\]
Assume there is a compact generating subcategory $\mathcal{G}\subset\mathcal{L}$
and natural equivalences
\[
J_{z,M,N}\colon F_0(M\otimes_zN)
\xrightarrow{\sim}
F_z(M)\widehat{\otimes}F_0(N)
\qquad (M,N\in\mathcal{G})
\]
compatible with units and the pentagon on triples of generators. Assume also
that natural endomorphisms of $F_0|_{\mathcal{G}}$ identify with $B$, whereas
natural endomorphisms of $F_z\widehat{\otimes}F_0$ identify with
$B\widehat{\otimes}B$.

Then there exists a unique meromorphic family of $A_\infty$-algebra maps
\[
\Delta_z\colon B\to B\widehat{\otimes}B
\]
transporting the action of $B$ across the $J_{z,M,N}$. The pentagon identity
forces the coassociativity relation
\[
(\mathrm{id}\widehat{\otimes}\Delta_w)\Delta_z
=
(\Delta_{z+w}\widehat{\otimes}\mathrm{id})\Delta_w.
\]
In particular, fusion can be reconstructed from one exact meromorphic fiber
functor together with generator-level coherence.
\end{theorem}

\begin{proof}
Transport the action of an element $b\in B$ on $F_0(M\otimes_zN)$ across
$J_{z,M,N}$. By the endomorphism-identification hypothesis, the transported
operator is given by a unique element $\Delta_z(b)\in B\widehat{\otimes}B$.
Naturality and multiplicativity of the original action imply that $\Delta_z$ is
an $A_\infty$-algebra map. Coassociativity is obtained by comparing the two
ways of transporting the action of $B$ across the two parenthesizations of
$F_0((L\otimes_zM)\otimes_wN)$; the pentagon says that the resulting
identifications agree, and hence so do the two induced coproducts.
\end{proof}

\begin{corollary}[Recognition on compact boundary modules \ClaimStatusProvedHere]
\label{cor:recognition-compact-boundary-modules}
Let $V_b\text{-}\mathrm{Mod}$ be a stable boundary category with vertex tensor
product $\boxtimes_z$, and let
\[
CB\colon V_b\text{-}\mathrm{Mod}\to \mathrm{Mod}^{\mathrm{comp}}(B)
\]
be derived conformal blocks. If a finite standard family compactly tensor-
generates $V_b\text{-}\mathrm{Mod}$, if $CB$ is fully faithful on standard tensor
words, and if generator-level fusion is intertwined with the meromorphic tensor
structure on $\mathrm{Mod}^{\mathrm{comp}}(B)$, then $CB$ is a monoidal
 equivalence on compact objects.
\end{corollary}

\begin{proof}
Apply Lemma~\ref{lem:generator-extension-exact-monoidal} to the compactly
generated subcategory spanned by the chosen standards.
\end{proof}

\section{Finite-stage RTT bar adjointness}
\label{sec:finite-stage-rtt-bar-adjointness}

Fix $V=\mathbf{C}^d$ and the rational Yang kernel
\[
R_\hbar(u)=1-\frac{\hbar P}{u}.
\]
For $N\ge 1$, let
\[
W_N:=\bigoplus_{r=1}^N \operatorname{End}(V)u^{-r}
\]
and write
\[
T(u)=1+\sum_{r=1}^N T^{(r)}u^{-r}.
\]
The finite RTT stage $A_N(R_\hbar)$ is the quadratic algebra generated by the
coefficients of the $T^{(r)}$ subject to coefficient extraction from the RTT
relation
\[
R_{12}(u-v)T_1(u)T_2(v)=T_2(v)T_1(u)R_{12}(u-v).
\]

\begin{lemma}[Scalar gauge invariance of finite RTT stages \ClaimStatusProvedHere]
\label{lem:scalar-gauge-invariance-finite-rtt}
If $R'(u)=f(u)R(u)$ with $f(u)\in 1+u^{-1}\mathbf{C}[[u^{-1}]]$ scalar and
invertible, then for every finite stage $N$ one has
\[
A_N(R')\cong A_N(R).
\]
In particular, the quadratic relation subspaces and the quadratic bar
corestrictions agree.
\end{lemma}

\begin{proof}
Because $f(u-v)$ is scalar,
\[
R'_{12}(u-v)T_1(u)T_2(v)-T_2(v)T_1(u)R'_{12}(u-v)
=
 f(u-v)\bigl(R_{12}(u-v)T_1(u)T_2(v)-T_2(v)T_1(u)R_{12}(u-v)\bigr).
\]
On the coefficient space truncated to modes $\le N$, multiplication by the
scalar series $f(u-v)$ is an invertible linear operator. Hence the spans of
the coefficient relations are identical.
\end{proof}

\begin{corollary}[Inverse kernel versus sign reversal \ClaimStatusProvedHere]
\label{cor:inverse-kernel-sign-reversal}
For the rational Yang kernel,
\[
R_\hbar(u)^{-1}
=
\frac{1}{1-\hbar^2/u^2}\left(1+\frac{\hbar P}{u}\right).
\]
Hence $R_\hbar^{-1}$ is scalar-gauge equivalent to $R_{-\hbar}$, and so
\[
A_N(R_\hbar^{-1})\cong A_N(R_{-\hbar}).
\]
\end{corollary}

\begin{definition}[Finite-stage bar corestriction and inverse-kernel cobracket]
\label{def:finite-stage-bar-corestriction}
Let $e_{ij}^{(r)}$ denote the normalized degree-$1$ bar generator dual to the
RTT generator $T_{ij}^{(r)}$. Define
\[
b_N^\hbar\colon (sW_N)^{\otimes 2}\to sW_N
\]
by
\[
b_N^\hbar([e_{ij}^{(r)}|e_{kl}^{(s)}])
=
-\delta_{kj}e_{il}^{(r+s-1)}+\delta_{il}e_{kj}^{(r+s-1)},
\]
with the convention that the right-hand side is zero when $r+s-1>N$.

Define also
\[
\delta_N^{-\hbar}\colon W_N^\vee\to W_N^\vee\otimes W_N^\vee
\]
by
\[
\delta_N^{-\hbar}(e_{ab}^{(m)})
=
\sum_{\substack{r,s\ge 1\\r+s-1=m}}\sum_{p=1}^d
\bigl(e_{ap}^{(r)}\otimes e_{pb}^{(s)}-e_{pb}^{(r)}\otimes e_{ap}^{(s)}\bigr).
\]
\end{definition}

\begin{theorem}[Finite-stage RTT bar adjointness \ClaimStatusProvedHere]
\label{thm:finite-stage-rtt-adjointness}
With respect to the coefficient pairing
\[
\langle e_{ij}^{(r)},e_{kl}^{(s)}\rangle
=\delta_{ik}\delta_{jl}\delta_{rs},
\]
one has
\[
b_N^\hbar=-(\delta_N^{-\hbar})^\vee.
\]
Equivalently, the finite-stage quadratic bar differential for $R_\hbar$ is the
strict adjoint of the inverse-kernel cobracket for
$R_\hbar^{-1}\simeq R_{-\hbar}$.
\end{theorem}

\begin{proof}
A direct computation gives
\[
\bigl\langle b_N^\hbar([e_{ij}^{(r)}|e_{kl}^{(s)}]),e_{ab}^{(m)}\bigr\rangle
=
\delta_{m,r+s-1}
\bigl(-\delta_{kj}\delta_{ia}\delta_{lb}
+\delta_{il}\delta_{ka}\delta_{jb}\bigr).
\]
On the other hand,
\[
\bigl\langle e_{ij}^{(r)}\otimes e_{kl}^{(s)},\delta_N^{-\hbar}(e_{ab}^{(m)})\bigr\rangle
=
\delta_{m,r+s-1}
\bigl(\delta_{kj}\delta_{ia}\delta_{lb}
-\delta_{il}\delta_{ka}\delta_{jb}\bigr).
\]
The two expressions differ by a minus sign, proving the claim.
\end{proof}

\begin{corollary}[The degree-$2$ twisting cochain \ClaimStatusProvedHere]
\label{cor:degree-two-twisting-cochain}
The adjoint map
\[
\tau_N:=-(b_N^\hbar)^\vee=\delta_N^{-\hbar}
\]
is the finite-stage degree-$2$ twisting cochain. Its Maurer--Cartan equation is
equivalent to the quadratic relation $(b_N^\hbar)^2=0$, hence to finite-stage
RTT/Yang--Baxter consistency.
\end{corollary}

\section{Corrected shifted RTT transport: bi-diagonal twist and boundary quotient}
\label{sec:corrected-shifted-rtt-transport}

\begin{remark}[Why diagonal conjugation is too weak]
\label{rem:diagonal-conjugation-too-weak}
Conjugating $T(u)$ by a diagonal matrix leaves the diagonal currents unchanged,
so it cannot encode the shifted-Yangian datum, whose extra input is precisely a
rational constraint on the Cartan currents. The right local surrogate is not
diagonal conjugation but a \emph{bi-diagonal left/right twist}. Even that twist
still lives at the local-kernel level only: it transports the quadratic RTT
calculus but does not yet produce the genuine shifted/truncated boundary data.
\end{remark}

\begin{definition}[Bi-diagonal twist]
\label{def:bi-diagonal-twist}
Choose rational series
\[
\ell_i^\mu(u),\ r_j^\mu(u)\in 1+u^{-1}\mathbf{k}[[u^{-1}]],
\qquad
\rho_{ij}^\mu(u):=\ell_i^\mu(u)r_j^\mu(u)
=\sum_{m\ge 0}\rho_{ij,m}^\mu u^{-m}.
\]
The associated lower-triangular automorphism
\[
\Phi_{\mu,N}\colon W_N\to W_N
\]
is defined by
\[
\Phi_{\mu,N}(t_{ij}^{(r)})
:=
\sum_{m=0}^{r-1}\rho_{ij,m}^\mu\,t_{ij}^{(r-m)}.
\]
Equivalently, the twisted currents are
\[
T_{ij}^{(\mu)}(u)=\rho_{ij}^\mu(u)T_{ij}(u).
\]
\end{definition}

\begin{theorem}[Transport of the finite RTT relation space \ClaimStatusProvedHere]
\label{thm:transport-of-finite-rtt-relation-space}
Let $\mathrm{Rel}_N$ denote the quadratic relation space of the unshifted finite
RTT stage, and let $\mathrm{Rel}_{\mu,N}$ denote the relation space extracted
from the twisted currents $T_{ij}^{(\mu)}(u)$. Then
\[
\mathrm{Rel}_{\mu,N}=(\Phi_{\mu,N}\otimes \Phi_{\mu,N})(\mathrm{Rel}_N).
\]
Hence the \emph{pre-shifted} finite stage
\[
A_{\mu,N}^{\mathrm{pre}}:=T(W_N)/\langle \mathrm{Rel}_{\mu,N}\rangle
\]
is canonically isomorphic to the unshifted stage.
\end{theorem}

\begin{proof}
Every coefficient of the RTT relation is a linear combination of quadratic
monomials $t_{ij}^{(r)}t_{kl}^{(s)}$ with coefficients coming from the fixed
kernel $R(u-v)$. Replacing $T_{ij}(u)$ by
$T_{ij}^{(\mu)}(u)=\rho_{ij}^\mu(u)T_{ij}(u)$ multiplies the coefficient of each
such monomial by the product $\rho_{ij}^\mu(u)\rho_{kl}^\mu(v)$. On the
coefficient space this is exactly the linear operator
$\Phi_{\mu,N}\otimes \Phi_{\mu,N}$. Therefore the relation space is transported
by that operator.
\end{proof}

\begin{theorem}[Finite-stage shifted bar adjointness \ClaimStatusProvedHere]
\label{thm:finite-stage-shifted-bar-adjointness}
Assume
\[
\rho_{ij}^{-\mu}(u)=\bigl(\rho_{ij}^{\mu}(u)\bigr)^{-1},
\qquad
\Phi_{-\mu,N}=\Phi_{\mu,N}^{-1}
\quad \text{mod }u^{-N-1}.
\]
Transport the coefficient pairing by
\[
\langle \Phi_{\mu,N}x,\Phi_{-\mu,N}y\rangle_\mu:=\langle x,y\rangle_0.
\]
Define transported structure maps by
\[
b_{N,\mu}
:=
\Phi_{\mu,N}b_N(\Phi_{\mu,N}^{-1}\otimes \Phi_{\mu,N}^{-1}),
\]
\[
\delta_{N,-\mu}
:=
(\Phi_{-\mu,N}\otimes \Phi_{-\mu,N})\,\delta_N\,\Phi_{-\mu,N}^{-1}.
\]
Then
\[
b_{N,\mu}=-(\delta_{N,-\mu})^\vee
\]
with respect to $\langle -, -\rangle_\mu$.
\end{theorem}

\begin{proof}
For $x,y,z\in W_N$,
\begin{align*}
\langle b_{N,\mu}(\Phi_{\mu,N}x\otimes \Phi_{\mu,N}y),\Phi_{-\mu,N}z\rangle_\mu
&=
\langle \Phi_{\mu,N}b_N(x\otimes y),\Phi_{-\mu,N}z\rangle_\mu \\
&=
\langle b_N(x\otimes y),z\rangle_0 \\
&=
-\langle x\otimes y,\delta_Nz\rangle_0 \\
&=
-\langle \Phi_{\mu,N}x\otimes \Phi_{\mu,N}y,
(\Phi_{-\mu,N}\otimes \Phi_{-\mu,N})\delta_Nz\rangle_\mu \\
&=
-\langle \Phi_{\mu,N}x\otimes \Phi_{\mu,N}y,
\delta_{N,-\mu}(\Phi_{-\mu,N}z)\rangle_\mu.
\end{align*}
This is exactly the stated adjointness.
\end{proof}

\begin{remark}[Twist-rigidity of the local Maurer--Cartan kernel]
\label{rem:twist-rigidity-local-mc}
If $\tau_N$ is the finite-stage degree-$2$ twisting cochain of
Corollary~\ref{cor:degree-two-twisting-cochain}, then its transported image
\[
\tau_{N,-\mu}:=\Phi_{-\mu,N}\tau_N\Phi_{\mu,N}^{-1}
\]
satisfies the Maurer--Cartan equation if and only if $\tau_N$ does. Thus the
local $r$-matrix/MC package is twist-rigid. The genuinely new shifted data do
not live here; they enter with the boundary quotient.
\end{remark}

\begin{definition}[Boundary quotient and orthogonal coideal]
\label{def:boundary-quotient-orthogonal-coideal}
Let $I_{\mu,\lambda}^{(N)}\subset A_{\mu,N}^{\mathrm{pre}}$ be the two-sided ideal
encoding the shifted/truncated boundary condition. Set
\[
A_{\mu,\lambda}^{(N)}:=A_{\mu,N}^{\mathrm{pre}}/I_{\mu,\lambda}^{(N)}.
\]
The associated \emph{orthogonal coideal} on the $-\mu$ side is
\[
C_{-\mu,\lambda}^{(N)}
:=
\operatorname{im}\bigl(\overline{B}(q_{\mu,\lambda}^{(N)})^\vee\bigr)
\subset \overline{B}(A_{-\mu,N}^{\mathrm{pre}}),
\]
where
\[
q_{\mu,\lambda}^{(N)}\colon A_{\mu,N}^{\mathrm{pre}}\to A_{\mu,\lambda}^{(N)}
\]
is the quotient map. Equivalently, $C_{-\mu,\lambda}^{(N)}$ is the annihilator of
$\ker \overline{B}(q_{\mu,\lambda}^{(N)})$ under the transported bar pairing.
\end{definition}

\begin{theorem}[Quotient/coideal descent \ClaimStatusProvedHere]
\label{thm:quotient-coideal-descent}
The orthogonal coideal $C_{-\mu,\lambda}^{(N)}$ is a dg subcoalgebra of
$\overline{B}(A_{-\mu,N}^{\mathrm{pre}})$, and there is a canonical
identification
\[
\overline{B}(A_{\mu,\lambda}^{(N)})^\vee\cong C_{-\mu,\lambda}^{(N)}
\]
weightwise. Hence the quotient bar differential on the $\mu$-side is adjoint to
the restricted cobar differential on the $-\mu$-side.
\end{theorem}

\begin{proof}
Functoriality of the reduced bar construction gives a dg coalgebra map
\[
\overline{B}(q_{\mu,\lambda}^{(N)})\colon
\overline{B}(A_{\mu,N}^{\mathrm{pre}})\to \overline{B}(A_{\mu,\lambda}^{(N)}).
\]
Its completed weightwise dual is a dg algebra map. The image of a dg algebra
map is stable under the dual differential, hence under the transported
coalgebraic differential. Under the perfect bar pairing supplied by
Theorem~\ref{thm:finite-stage-shifted-bar-adjointness}, this image is exactly the
annihilator of the quotient kernel. Since the dual of a surjective map is
injective, the image identifies canonically with the dual of the quotient.
\end{proof}

\[
A_{\mu,N}^{\mathrm{pre}}\twoheadrightarrow A_{\mu,\lambda}^{(N)},
\qquad
C_{-\mu,\lambda}^{(N)}\hookrightarrow \overline{B}(A_{-\mu,N}^{\mathrm{pre}}),
\qquad
\overline{B}(A_{\mu,\lambda}^{(N)})^\vee\cong C_{-\mu,\lambda}^{(N)}.
\]

\section{The rank-one test case: from orthogonal coideals to Kleinian surfaces}
\label{sec:rank-one-test-case}

The first nontrivial case is already instructive enough to correct the overall
programme.

\subsection{Why the naive rank-one Cartan quotient fails}

\begin{proposition}[Stage $N=1$ twist is trivial; naive Cartan fixing collapses the root sector \ClaimStatusProvedHere]
\label{prop:stage-one-cartan-collapse-appendix}
For $A_1$ at stage $N=1$, the bi-diagonal twist is trivial on generators:
\[
\Phi_{\mu,1}=\mathrm{id}.
\]
Any quotient of $U(\mathfrak{sl}_2)$ in which $h$ becomes a scalar
collapses the root sector: the images of both $e$ and $f$ vanish.
\end{proposition}

\begin{proof}
At $N=1$ one has
\[
\Phi_{\mu,1}(t_{ij}^{(1)})=\rho_{ij,0}^\mu t_{ij}^{(1)}=t_{ij}^{(1)},
\]
so the twist is the identity.

For the second statement, suppose $h$ maps to a scalar $\beta$. Then in the
quotient,
\[
0=[\beta,e]=[h,e]=2e,
\qquad
0=[\beta,f]=[h,f]=-2f.
\]
Hence $e=f=0$.
\end{proof}

\begin{remark}[The first non-collapsing boundary quotient]
\label{rem:first-non-collapsing-boundary-quotient}
Proposition~\ref{prop:stage-one-cartan-collapse} shows that the first genuinely
nontrivial rank-one boundary condition cannot be a linear scalar fixing of the
Cartan mode. The minimal non-collapsing quotient is a \emph{central quadratic}
one: in RTT language this is the first determinantal/Casimir condition. This
is the correct rank-one shadow of the shifted/truncated boundary layer.
\end{remark}

\subsection{The $A_1$ determinantal/Casimir boundary quotient}

Let $e,f,h$ be the standard generators of $U(\mathfrak{sl}_2)$ with
\[
[h,e]=2e,
\qquad
[h,f]=-2f,
\qquad
[e,f]=h.
\]
Introduce the Casimir-type central element
\[
\Omega:=4fe+h^2+2h=4ef+h^2-2h.
\]
For $c\in\mathbf{k}$, define the central quotient
\[
A_c:=U(\mathfrak{sl}_2)/(\Omega-c).
\]
This is the first nontrivial rank-one boundary quotient that preserves the root
sector.

\begin{definition}[Rank-one orthogonal coideal]
\label{def:rank-one-orthogonal-coideal}
Let
\[
q_c\colon U(\mathfrak{sl}_2)\to A_c
\]
be the quotient map. The associated orthogonal coideal is
\[
C_c:=\operatorname{im}\bigl(\overline{B}(q_c)^\vee\bigr)
\subset \overline{B}(U(\mathfrak{sl}_2)).
\]
Equivalently, $C_c$ is the annihilator of the bar ideal generated by the curved
relation
\[
\Omega-c=4fe+h^2+2h-c.
\]
\end{definition}

\begin{theorem}[Explicit rank-one orthogonal coideal and curved cobar model \ClaimStatusProvedHere]
\label{thm:explicit-rank-one-orthogonal-coideal}
The orthogonal coideal $C_c$ is the curved dg subcoalgebra of
$\overline{B}(U(\mathfrak{sl}_2))$ obtained by deleting the one-dimensional bar
direction dual to the Casimir relation
\[
4f|e+h|h+2h-c.
\]
Its curved cobar algebra is
\[
\Omega_\theta(C_c)
\cong
T(e,f,h)\Big/\bigl(he-eh-2e,\ hf-fh+2f,\ ef-fe-h,\ 4fe+h^2+2h-c\bigr)
\cong A_c.
\]
\end{theorem}

\begin{proof}
The bar construction of the quotient map $q_c$ is a morphism of curved dg
coalgebras. By definition, $C_c$ is the image of the dual map, hence the
annihilator of the coideal generated by the extra relation $\Omega-c$. The
linear and quadratic pieces of that relation are the bar words
\[
2h,
\qquad
4f|e+h|h,
\]
while the constant term contributes the curvature parameter $c$.
Consequently the cobar algebra of $C_c$ is the free tensor algebra on $e,f,h$
with the usual $\mathfrak{sl}_2$ relations and the additional relation
$4fe+h^2+2h-c=0$. That is exactly the quotient $A_c$.
\end{proof}

\begin{corollary}[Generalized Weyl normal form \ClaimStatusProvedHere]
\label{cor:generalized-weyl-normal-form}
Let $\sigma\in \operatorname{Aut}(\mathbf{k}[h])$ be given by
\[
\sigma(h)=h-2,
\]
and define the polynomial
\[
a_c(h):=\frac{1}{4}\bigl(c-h(h+2)\bigr).
\]
Then
\[
A_c\cong \mathbf{k}[h](\sigma,a_c),
\]
the rank-one generalized Weyl algebra generated by $x,y,h$ with
\[
xy=\sigma(a_c),
\qquad yx=a_c,
\qquad xr=\sigma(r)x,
\qquad yr=\sigma^{-1}(r)y
\quad (r\in\mathbf{k}[h]).
\]
Under this identification one may take
\[
x=e,
\qquad y=f.
\]
\end{corollary}

\begin{proof}
The commutation relations imply
\[
eh=(h-2)e,
\qquad fh=(h+2)f,
\]
so $e$ and $f$ implement the automorphism $\sigma$ and its inverse. The Casimir
relation gives
\[
fe=\frac{1}{4}\bigl(c-h(h+2)\bigr)=a_c(h),
\]
while the identity $\Omega=4ef+h^2-2h$ yields
\[
ef=\frac{1}{4}\bigl(c-h(h-2)\bigr)=a_c(h-2)=\sigma(a_c)(h).
\]
These are precisely the defining relations of the generalized Weyl algebra
$\mathbf{k}[h](\sigma,a_c)$.
\end{proof}

\begin{corollary}[Kleinian associated graded at the nilpotent point \ClaimStatusProvedHere]
\label{cor:kleinian-associated-graded}
Equip $A_c$ with the PBW filtration in which $e,f,h$ have degree $1$. At the
nilpotent point $c=0$ one has
\[
\operatorname{gr}A_0
\cong
\mathbf{k}[e,f,h]/(4ef+h^2).
\]
After a linear change of variables this is the coordinate ring of the affine
Kleinian surface of type $A_1$:
\[
\mathbf{k}[x,y,z]/(xy-z^2).
\]
\end{corollary}

\begin{proof}
In the associated graded algebra, the commutator relations become commutative,
whereas the leading term of the extra relation $4fe+h^2+2h=0$ is
$4ef+h^2=0$. Thus
\[
\operatorname{gr}A_0\cong \mathbf{k}[e,f,h]/(4ef+h^2).
\]
Set $x=2e$, $y=2f$, and $z=\sqrt{-1}\,h$ over $\mathbf{C}$; then the relation
becomes $xy-z^2=0$.
\end{proof}

\begin{remark}[Interpretation]
\label{rem:rank-one-interpretation}
The rank-one orthogonal coideal already captures the local algebraic music of
the boundary layer. The local RTT kernel is unchanged at $N=1$;
nontriviality enters by a central quadratic boundary quotient; the dual object is
a curved orthogonal coideal; and its cobar algebra is the quantized rank-one
Kleinian surface. This is the first concrete model of the principle that
``shifted/truncated data live in the boundary quotient, not in the local kernel
alone''.
\end{remark}

\begin{conjecture}[Rank-one BFN matching problem \ClaimStatusConjectured]
\label{conj:rank-one-bfn-matching}
After translating parameters between the Yangian/Casimir normalization and the
Coulomb-branch normalization, the rank-one orthogonal coideal $C_c$ should agree
with the finite-stage $A_1$ instance of the general orthogonal coideal
$C_{-\mu,\lambda}^{(N)}$, and the corresponding cobar algebra should coincide
with the quantized Coulomb branch of the rank-one $A_1$ theory.

Equivalently: the remaining task in rank one is not bar/cobar homological
algebra but the explicit identification of the boundary quotient parameter $c$
with the Coulomb-branch mass/FI data.
\end{conjecture}

\section{Synthesis and the fortified frontier}
\label{sec:synthesis-fortified-frontier}

\begin{theorem}[Fortified frontier package \ClaimStatusProvedHere]
\label{thm:fortified-frontier-package}
Within the ordered Yangian programme, the following statements are now formal.
\begin{enumerate}[label=\textup{(\arabic*)}]
\item Positive weight towers admit a stabilized completed bar/cobar envelope and
 weightwise Maurer--Cartan descent
 \textup{(}Theorems~\ref{thm:stabilized-completed-recovery}
 and~\ref{thm:weightwise-maurer-cartan-descent}\textup{)}.
\item A finite bar-duality datum yields an explicit chain-level duality map
 $\Theta_A$ and hence a tangent-level Lagrangian envelope
 \textup{(}Theorem~\ref{thm:trace-theoretic-bar-duality} and
 Corollary~\ref{cor:lagrangian-envelope}\textup{)}.
\item The local finite RTT kernel is self-dual under
 $R\mapsto R^{-1}$ by strict bar/cobar adjointness
 \textup{(}Theorem~\ref{thm:finite-stage-rtt-adjointness}\textup{)}.
\item The correct shifted finite-stage transport is a bi-diagonal twist; it
 preserves the local Maurer--Cartan package and transports bar adjointness
 exactly \textup{(}Theorems~\ref{thm:transport-of-finite-rtt-relation-space}
 and~\ref{thm:finite-stage-shifted-bar-adjointness}\textup{)}.
\item The genuinely new shifted/truncated datum lives in the boundary quotient,
 whose dual avatar is the orthogonal coideal
 \textup{(}Theorem~\ref{thm:quotient-coideal-descent}\textup{)}.
\item In the first nontrivial rank-one case, that orthogonal-coideal mechanism
 already produces a curved cobar model whose associated graded is the affine
 Kleinian surface of type $A_1$
 \textup{(}Theorem~\ref{thm:explicit-rank-one-orthogonal-coideal} and
 Corollary~\ref{cor:kleinian-associated-graded}\textup{)}.
\end{enumerate}
\end{theorem}

\begin{proof}
Each clause is the content of the cited theorem or corollary. No new computation is needed; the programme has been split
into the right pieces: local kernel, completed envelope, boundary quotient,
orthogonal dual, and rank-one test case.
\end{proof}

\begin{remark}[What remains open]
\label{rem:what-remains-open}
The remaining work is now sharply localized.
\begin{enumerate}[label=\textup{(\roman*)}]
\item Construct the genuine boundary quotient $I_{\mu,\lambda}^{(N)}$ in a way that
 is manifestly compatible with the shifted Yangian and BFN presentations.
\item Identify the orthogonal coideal $C_{-\mu,\lambda}^{(N)}$ with the expected
 dg-shifted or Coulomb-branch dual presentation.
\item Lift the rank-one Casimir/Kleinian model to higher finite stages and then
 to completed towers.
\item Match the reconstructed meromorphic coproduct from conformal blocks with
 the line-operator coproduct and $r(z)$ of \cite{DNP25} on compact generators.
\item Globalize the tangent-level Lagrangian envelope to a genuine shifted-
 symplectic deformation theory on the factorization side.
\end{enumerate}
The major change is that none of these are now hidden inside a blurred slogan.
They are separate and exact problems.
\end{remark}

\section{Conventions and bridges}
\label{sec:shifted-rtt-conventions}

\subsection{Three layers and the four Yangian types}
\label{subsec:shifted-rtt-yangian-types}

\begin{convention}[Layers (a)/(b)/(c) match Yangian types
1/2/$+$shift; \ClaimStatusProvedHere]
\label{conv:shifted-rtt-yangian-types}
The three-layer decomposition of
Definition~\ref{def:three-layers-ordered-theory-appendix} aligns
with the four Yangian types as follows:
\begin{enumerate}[label=\textup{(\arabic*)}]
\item Layer (a) (\emph{local RTT kernel}): \emph{type 1},
the classical Yangian $Y_\hbar(\mathrm{sl}_N)$ as
$E_1$-topological algebra on $\R$, RTT presentation.
\item Layer (b) (\emph{completed envelope}): \emph{type 2},
the dg-shifted Yangian $Y^{\mathrm{dg}}_\hbar(\mathrm{sl}_N)$ at
the formal disk, with explicit $A_\infty$ structure on the
weightwise stabilization (Lemma~\ref{lem:weightwise-stabilization-bar-complex}, Definition~\ref{def:stabilized-completed-bar-coalgebra}).
\item Layer (c) (\emph{boundary quotient}): the shifted/truncated/
Coulomb-branch datum, which is genuinely additional to the
four-Yangian-types axis and corresponds to the coweight $\mu$
in the unified $Q_\fg^{k,f,\mu}$ of
Theorem~\ref{thm:unified-chiral-quantum-group}.
\end{enumerate}
The chiral Yangian $Y(\mathrm{sl}_N)^{\mathrm{ch}}$ (type 3) and the
spectral factorization on $\mathbb{A}^1_u$ (type 4) appear via the
factorization-homology lift of layer (a)+(b) along the curve $X$
(Vol~I~\cref{V1-thm:theorem-A-E1}).
\end{convention}

\subsection{Working notion: shifted RTT in Notion~B}
\label{subsec:shifted-rtt-notion-b}

\begin{convention}[Shifted RTT data are Notion~B
$E_1$-chiral; \ClaimStatusProvedElsewhere]
\label{conv:shifted-rtt-notion-b}
The shifted RTT presentation, the orthogonal coideal mechanism,
and the rank-one Kleinian test case are all Notion-B
($A_\infty$ in $\mathrm{End}^{\mathrm{ch}}$): structure maps
chain-level, Stasheff identities holding, the bar coalgebra is
the deconcatenation $T^c$ on the desuspended augmentation ideal.
The strictification beyond the local quadratic kernel
corresponds to nontrivial $m_k$ for $k\ge 3$.
\end{convention}

\subsection{Coideal property: chain vs differential stability}
\label{subsec:shifted-rtt-coideal-discipline}

\begin{remark}[Coideal vs sub-dg-module distinction;
\ClaimStatusProvedElsewhere]
\label{rem:shifted-rtt-coideal-discipline}
``Image of a dg algebra map is stable under the dual differential''
is differential stability, NOT coproduct stability (the actual
coideal property $\Delta(C) \subset C\otimes\text{ambient} +
\text{ambient}\otimes C$). Theorem~\ref{thm:quotient-coideal-descent}
is read here in its strongest honest form: the descent claim
yields a sub-dg-module that becomes a coideal upon the explicit
verification of the $\Delta$-stability, which holds in the
rank-one Kleinian test case
(Theorem~\ref{thm:explicit-rank-one-orthogonal-coideal}) and is
verified case-by-case at higher rank as the construction is
lifted to the completed tower.
\end{remark}

\subsection{Terminology: ``orthogonal coideal'' vs Letzter--Kolb}
\label{subsec:shifted-rtt-terminology}

\begin{remark}[Disambiguation: ``orthogonal'' here means
\emph{pairing-annihilator}; \ClaimStatusProvedElsewhere]
\label{rem:shifted-rtt-terminology}
The term ``orthogonal coideal'' in the
shifted-Yangian/B-type literature (Letzter, Kolb) refers to
$\imath$-quantum groups arising from symmetric pairs (pairs
$(\mathfrak{g}, \mathfrak{g}^\theta)$ with $\theta$ an
involution); it is a representation-theoretic structure on
quantum groups associated to symmetric pairs $\mathrm{B}$/$\mathrm{D}$.
In this chapter, ``orthogonal coideal'' refers exclusively to
the bar-pairing-annihilator (Definition above), an entirely
different structure: a coideal in a chiral Hopf algebra defined as
the orthogonal complement of a subalgebra under the bar pairing.
The two structures live in different categories and should not
be conflated. Cross-volume readers using Letzter--Kolb terminology
should apply the renaming ``pairing-annihilator coideal'' to
disambiguate.
\end{remark}

\subsection{Level-prefix discipline for the local RTT kernel}
\label{subsec:shifted-rtt-rmatrix-level-prefix}

\begin{remark}[Level prefix $\hbar$ on the RTT kernel;
\ClaimStatusProvedElsewhere]
\label{rem:shifted-rtt-rmatrix-level-prefix}
The local RTT kernel $R(u, v) = 1 + \hbar P/(u-v) + O(\hbar^2)$
on $V\otimes V$ ($V = \C^N$, $P$ permutation, $u-v$ spectral
displacement) carries the level prefix $\hbar$. Substituting
$\hbar = 0$ gives $R(u,v) = 1$ (trivial $R$-matrix, abelian limit):
substituting the level parameter to zero gives the trivial
$r$-matrix on the trace-form side. The shifted $R$-matrix
$R_\mu(u, v) = R(u, v) \cdot
\prod_{i=1}^{N-1} (u-v+\mu_i\hbar)$
introduces the shift datum $\mu$ as additional structure in
layer (c); the level prefix $\hbar$ is preserved.
\end{remark}

\subsection{Shifted RTT as $\mu \ne 0$ specialisation
of the unified $Q_\fg^{k,f,\mu}$}
\label{subsec:shifted-rtt-holography-anchor}

\begin{remark}[Shifted RTT as the $\mu \ne 0$ fibre of the unified
$Q_\fg^{k,f,\mu}$; \ClaimStatusProvedElsewhere]
\label{rem:shifted-rtt-uhm}
The shifted RTT and orthogonal coideal package of this appendix
is the $\mu \ne 0$ specialisation of the unified chiral quantum
group $Q_\fg^{k,f,\mu}$
(Theorem~\ref{thm:unified-chiral-quantum-group}) at
$(\fg, k, f) = (\mathrm{sl}_N, k\text{-generic}, 0)$ with
arbitrary shift coweight $\mu \in P(\mathrm{sl}_N)^+$. The
weightwise stabilization mechanism of layer (b) provides the
chain-level $A_\infty$ control needed to upgrade the Yangian
fibre to the shifted variant. The strict factorization quantum
group output of \S\ref{chap:dg-shifted-factorization} extends
to the shifted variant by the same root-multiplicity-one mechanism,
modulo the verification of layer (c) data; this is the
type-A specialisation of the universal holography input
to \cref{thm:universal-holography-master} (Vol~II
\cref{ch:programme-climax-platonic}). The Vol~I unified Koszul reflection
(\cref{V1-thm:koszul-reflection}) controls the bar-side duality;
the universal conductor (\cref{V1-chap:universal-conductor})
locates the shifted variant on the conductor strata via the
shift coweight $\mu$.
\end{remark>
